\documentclass[9pt, aspectratio=169]{beamer}
\usepackage{FiraSans}
\usetheme[subsectionpage=progressbar]{metropolis}
\usepackage[utf8]{inputenc}
\usepackage{amsmath}
\usepackage{amsfonts}
\usepackage{amssymb}
\usepackage{multicol}
\usepackage{tikz}
\usepackage{caption}
\usepackage{xcolor}
\usepackage[absolute,overlay]{textpos}
\usepackage[T1]{fontenc} 
\usepackage[skins]{tcolorbox}
\author{Nicola Roman\`o - nicola.romano@ed.ac.uk}
\title{Lecture 20 - Recents advances in image analysis using deep learning\\\small Part 2 - Image segmentation}
\setlength{\fboxsep}{0pt}
\setbeamertemplate {footline}{\begin{scriptsize}\hfill\insertframenumber ~of \inserttotalframenumber\kern1em\vskip5pt\end{scriptsize}}

% Remove "Figure" in front of captions
% See https://tex.stackexchange.com/questions/82456/how-to-remove-figure-caption-prefix-figure-in-beamer
\captionsetup{labelformat=empty,labelsep=none}

\titlegraphic{\centering \includegraphics[scale=.5]{instituteLogo.png}}
\date{}

\begin{document}

\newtcolorbox{codebox}{enhanced,
    top=2pt,
    left=2pt,
    right=2pt,
    bottom=2pt,
    boxrule=0pt,
    leftrule=5pt,
    sharp corners,
    colback=gray!20,
    colframe=blue!60!black}

\begin{frame}
    \titlepage
\end{frame}

\begin{frame}
    {Learning objectives}
    \begin{columns}
        \begin{column}{0.8\textwidth}
            \begin{itemize}
                \item Discuss recent advances in image analysis using deep learning.
            \end{itemize}

            \pause
            \vspace{2em}

            Today we are going to analyse recent articles related to image segmentation using deep learning.

            \vspace{1em}
            \begin{itemize}
                \item The StarDist and StarDist3D methods for cell segmentation.
                      \begin{itemize}
                          \item Schmidt et al. 2018 - Cell Detection with Star-convex Polygons
                          \item Weigert et al. 2020 - Star-convex Polyhedra for 3D Object Detection and Segmentation in Microscopy
                      \end{itemize}
                \item TissueNet and Mesmer
                      \begin{itemize}
                          \item Greenwald et al. 2021 - Whole-cell segmentation of tissue images with human-level performance using large-scale data annotation and deep learning
                      \end{itemize}
            \end{itemize}
        \end{column}
        \begin{column}{0.2\textwidth}
            \includegraphics[angle=-30, origin=tr, width=1.5\textwidth]{lightbulb.png}
        \end{column}
    \end{columns}
\end{frame}

\section{StarDist and StarDist3D}

\begin{frame}
    {Instance segmentation of cell nuclei}
    Instance segmentation of microscopy images is a difficult task.\\

    \textit{Top-down approaches} perform semantic segmentation, then split instances.\\
    \textbf{Problem}: nearby cells in crowded images fused together.\\

    \textit{Bottom-up approaches} roughly define cells (bounding boxes), then classify pixels within each box.\\
    \textbf{Problem}: NMS steps might miss nearby nuclei.

    \centering
    \includegraphics[width=.5\textwidth]{Schmidt2018_problems.png}
\end{frame}

\begin{frame}
    {StarDist: cell detection with star-convex polygons}
    \centering
    \includegraphics[width=\textwidth]{Schmidt2018_title.png}

    \raggedright
    \textit{``To alleviate the aforementioned problems, we propose StarDist, a cell detection method that predicts a shape representation which is flexible enough such that – without refinement – the accuracy of the localization can compete
        with that of instance segmentation methods. To that end, we use \textbf{star-convex polygons} that we find well-suited to approximate the typically roundish shapes of cell nuclei in microscopy images.''}
\end{frame}

\begin{frame}
    {StarDist: method}
    \centering
    \includegraphics[width=.5\textwidth]{Schmidt2018_process.png}

    \begin{textblock*}{2cm}(2cm,2.5cm)
        Radial distances to border
    \end{textblock*}

    \begin{textblock*}{2cm}(12cm,2.5cm)
        Object probabilities
    \end{textblock*}

    \begin{textblock*}{2cm}(2cm,6.5cm)
        U-Net architecture
    \end{textblock*}

    \begin{textblock*}{4cm}(12cm,7cm)
        Obj. prob. $\rightarrow$ binary cross-entropy loss
    \end{textblock*}

    \begin{textblock*}{4cm}(12cm,6cm)
        Distances $\rightarrow$ MAE loss
    \end{textblock*}
\end{frame}

\begin{frame}
    {Results}
    \centering
    \only<1>{
        \includegraphics[width=.75\textwidth]{Schmidt2018_comparison.png}
    }
    \only<2>{
        \includegraphics[width=\textwidth]{Schmidt2018_comparison2.png}
    }

    \footnotesize
    \raggedright
    StarDist outperforms many common instance segmentation methods.
\end{frame}

\begin{frame}
    {Can we extend StarDist to 3D?}
    \centering
    \includegraphics[width=\textwidth]{Weigert2020_title.png}
\end{frame}

\begin{frame}
    {StarDist3D - method}
    \centering
    \includegraphics[width=\textwidth]{Weigert2020_method.png}
\end{frame}

\begin{frame}
    {Results}
    \centering

    \includegraphics[width=\textwidth]{Weigert2020_results.png}
\end{frame}

\begin{frame}
    {Comparison with existing methods}
    \centering
    \includegraphics[width=\textwidth]{Weigert2020_comparison.png}
\end{frame}

\section{TissueNet}

\begin{frame}
    {The problem}
    \begin{itemize}
        \item Lack of a generalized algorithm for locating single cells in images
        \item Deep learning requires a lot of ground-truth data (difficult to get!) for high accuracy segmentation
        \item Public datasets only annotate nuclei, not the entire cell
    \end{itemize}
\end{frame}
\begin{frame}
    {TissueNet and the Mesmer algorithm}
    \centering
    \includegraphics[width=\textwidth]{Greenwald2021_title.png}

    \begin{itemize}
        \item \textbf{TissueNet}: ``a comprehensive segmentation dataset of >1 million paired whole-cell and nuclear annotations.''
        \item \textbf{Mesmer}: a deep-learning segmentation algorithm.
        \item \textbf{DeepCell}: an open-source collection of software libraries, to create a web interface for using Mesmer, as well as plugins for ImageJ and QuPath.
    \end{itemize}
\end{frame}

\begin{frame}
    {TissueNet development}
    \centering
    \includegraphics[width=\textwidth]{Greenwald2021_annotation_process.png}

    \footnotesize
    \raggedright
    TissueNet was developed through a crowdsourcing approach
\end{frame}

\begin{frame}
    {TissueNet - some statistics}
    \centering
    \includegraphics[width=\textwidth]{Greenwald2021_stats.png}
\end{frame}

\begin{frame}
    {The Mesmer segmentation algorithm}
    \centering
    \includegraphics[width=.8\textwidth]{Greenwald2021_Mesmer_method.png}
\end{frame}

\begin{frame}
    {Mesmer architecture}
    \centering
    \includegraphics[width=\textwidth]{Greenwald2021_architecture.png}
\end{frame}

\begin{frame}
    {Segmentation results}
    \centering
    \only<1>{
        \includegraphics[width=.9\textwidth]{Greenwald2021_method_comparison.png}
    }
    \only<2>{
        \includegraphics[width=.5\textwidth]{Greenwald2021_method_F1_score.png}
    }
\end{frame}

\begin{frame}
    {Segmentation of different tissues}
    \centering
    \includegraphics[width=\textwidth]{Greenwald2021_various_tissues.png}
\end{frame}

\begin{frame}
    {Human-level performance in segmentation}
    \centering
    \includegraphics[width=.7\textwidth]{Greenwald2021_Mesmer_vs_human.png}
\end{frame}

\begin{frame}
    {Mesmer deployment}
    \centering
    \includegraphics[width=.75\textwidth]{Greenwald2021_Mesmer_deployment.png}
\end{frame}

\begin{frame}
{The end!}
\centering
\Large
That's the end of the course\\
Hope you enjoyed it!

\footnotesize
(and looking forward to see your assignments!)
\end{frame}

\end{document}